


\section{Memory Experiment Configuration and Architectural Disclosure}
\label{sec:appendix_memory}



This appendix documents the memory configurations used in the MemoryAgentBench experiments and discloses the architectural constraints under which each evaluated framework operates. The intent of this appendix is not to claim full equivalence or strict normalization across memory implementations, but rather to transparently describe how memory was instantiated, controlled, and isolated in practice, and to explain why several observed behaviors arise from architectural design choices rather than parameter tuning or experimental artifacts. As such, this appendix should be read as an architectural disclosure that contextualizes the results reported in the main paper.

\textbf{Common Experimental Configuration.}
All memory experiments follow a shared configuration to ensure comparability across frameworks. Each agent is evaluated using the same base language model (Open AI GPT-OSS:20b) with identical decoding parameters, including a maximum generation length of 1500 tokens and a low temperature of 0.1 to minimize stochastic variation. Context ingestion is performed using fixed-size overlapping chunks, with a maximum chunk size of 4096 tokens and an overlap of 200 tokens to preserve semantic continuity across chunk boundaries. The total context budget is explicitly bounded and swept across experiments to study scaling behavior and memory saturation effects. All MemoryAgentBench splits are evaluated, including Accurate Retrieval, Test-Time Learning, Long-Range Understanding, and Conflict Resolution. Each benchmark run is executed in a fully isolated session. For frameworks with persistent memory backends, a fresh storage directory or database file is created per session and removed after evaluation, preventing cross-contamination between runs and ensuring that memory behavior reflects only the current benchmark instance.

\textbf{Memory Interfaces and Normalization Strategy.}
To enable systematic evaluation across heterogeneous frameworks, all systems are wrapped behind a common adapter interface exposing three operations: \texttt{reset()}, \texttt{ingest(context)}, and \texttt{query(question)}. This interface standardizes how benchmark context is provided and how questions are issued, but it does not modify or emulate internal memory semantics. Consequently, memory behavior reflects native architectural properties rather than synthetic normalization. Importantly, no evaluated framework supports explicit memory editing, deletion, or overwrite operations. All memory backends are append-only, whether implemented through vector stores, SQLite persistence, or conversational accumulation. As a result, selective forgetting is evaluated indirectly through retrieval behavior, context truncation, or recency bias rather than through explicit memory manipulation.

\textbf{Framework-Specific Memory Realizations.}
The OpenAI Agents SDK implements memory through a persistent SQLite-backed session that stores accumulated interaction history. Long-term memory is realized as conversation accumulation, while short-term memory is governed by an explicit context window limit enforced through input truncation. During ingestion, short-term accumulation is disabled by resetting the session after each chunk is stored, ensuring that ingestion does not inflate prompt context. During querying, accumulated memory is truncated according to a configurable context budget that is swept across experiments. As a result, memory behavior in the SDK is dominated by context budgeting rather than retrieval selectivity, and performance improves with larger context windows at the cost of increased latency and token usage.

CrewAI provides built-in short-term, long-term, and entity memory abstractions backed by persistent storage. In our experiments, memory is enabled at the framework level and persists across agent interactions within a session. Context ingestion is performed by a dedicated ingest agent that stores chunked benchmark content into CrewAI’s memory backend, while querying is handled by a separate answering agent that relies on the framework’s internal retrieval mechanisms. Memory resets are implemented by instantiating a new CrewAI environment with a fresh storage directory. Although CrewAI exposes multiple memory categories, retrieval behavior remains opaque and does not provide explicit controls for relevance thresholds, ordering, or deletion.


Agno implements memory as a SQLite-backed store that automatically captures user-provided content and injects it into the prompt at query time, resulting in accumulation-based memory behavior within each session. In our experiments, memory ingestion is performed by issuing explicit memorization commands for each context chunk, and memory resets are implemented by creating a fresh database per session. Under this configuration, Agno does not expose explicit controls for retrieval limits, relevance thresholds, or selective deletion, leading to unfiltered accumulation of stored content during a session. We note that recent versions of Agno introduce a documented memory optimization mechanism based on LLM-driven summarization and compression, which reduces token usage by collapsing multiple memories into fewer summarized entries. While this capability can significantly reduce context size, it fundamentally alters memory granularity and retrieval semantics by replacing fine-grained entries with aggregated summaries. To preserve comparability with other accumulation-based frameworks and to avoid introducing an additional summarization step that is not uniformly available across systems, this optimization was not enabled in our evaluation. As a result, the reported Agno results reflect native accumulation behavior without post hoc memory compression. In contrast, the effect of increasing context window size on accumulation-based memory performance is explicitly studied in the OpenAI Agents SDK experiments, where context budgeting is directly controlled and swept as a primary experimental variable. 


LangGraph is evaluated in two configurations. In the retrieval-only configuration, LangGraph uses an in-memory semantic store backed by embeddings to represent long-term memory. Context is ingested as chunked documents stored in a vector index, and at query time a fixed number of relevant chunks are retrieved and injected into a stateless prompt. No short-term message accumulation is used, and each query is independent. This configuration isolates semantic retrieval behavior and avoids context overflow, but it cannot support test-time learning or incremental reasoning across turns. In the hybrid configuration, LangGraph combines semantic long-term retrieval with short-term message accumulation via a checkpointer. Retrieved memories are injected alongside a truncated conversation history bounded by a token budget. This design mirrors accumulation-based memory while retaining retrieval capabilities. However, because short-term memory is append-only and truncation is purely recency-based, contradictory or outdated information cannot be removed, which directly impacts selective forgetting and conflict resolution.

\textbf{Implications for Interpretation.}
Taken together, these configurations demonstrate that memory behavior in current multi-agent frameworks is fundamentally constrained by architectural design rather than tunable parameters. Retrieval-centric systems excel at accurate recall and bounded context usage but struggle to adapt over time. Accumulation-based systems support test-time learning but incur high token costs and degrade under long horizons. Hybrid systems inherit both advantages and failure modes, resulting in fragile behavior when memory grows or conflicts arise. Accordingly, the results reported in the main paper should be interpreted as properties of memory paradigms rather than implementation artifacts. No evaluated framework provides native support for explicit memory revision, contradiction resolution, or selective deletion, which explains the uniformly poor selective forgetting performance observed across systems.
